\cotefile{157-5}
\scanpage{389}\mspage{303}

% slanted
\margin{Refondre avec 4.13.23 etc.}

\uline{Définition} 4.13.11. Soient $M$ une catégorie \deleted{$\pi_{\unsure{0}}$}-accessible, $\pi_0$ un cardinal utile pour $M$, $\pi\inserted{\unsure{>|\geq}\pi_0}$ un cardinal. \deleted{\unsure{infini}} Une catégorie $I$ est dite $\pi_0$-$\pi$-\uline{adaptée} si a) $I$ est équivalente à une petite catégorie, b) $I$ est \unsure{grde|grd} devant $\pi_0$, et c) $I$ est stable par $\underset{\longrightarrow}{\lim}$ de cardinal $\leq\pi$, et \unsure{grdes|grde} devant $\pi_0$.

% slanted
\margin{\uline{NB} Si $\pi=\pi_0$, la condition c) équivaut à dire que $I$ est stable par facteurs directs.}

Un système inductif $\inserted{(x_i)_{i\in I}}$ de type $I$ dans $M$ est dit $\pi_0$-$\pi$-adapté, si $I$ est $\pi_0$-$\pi$-adapté i.e. satisfait les conditions a) b) c), et si de plus on a d) le foncteur $i\longrightarrow x_i:I\longrightarrow M$ commute aux $\underset{\longrightarrow}{\lim}$ (\unsure{grdes|grde} devant $\pi_0$ et de cardinal $\leq\pi$) envisagées dans c), et e) les $x_i$ sont dans $M_\pi$.

\uline{Proposition} 4.13.12. Soient $M$ une catégorie $\pi_0$-accessible, et $\pi>\pi_0$ un cardinal utile pour $M$. Alors tout $x\in\mathrm{Ob}M$ admet une représentation comme $\deleted{x\simeq}\underset{j\in J}{\underset{\longrightarrow}{\lim}}\,x_j$, où le système inductif $(x_j)_{j\in J}$ dans $M$ est $\pi_0$-$\pi$-adapté. De façon précise, on peut prendre $\deleted{\unsure{I}}J=M_\pi/x$.

% archive sheet number: 388
\scanpage{390}\mspage{304}

En effet\unsure{,|.} la vérification des propriétés a) à d), pour le système inductif canonique, de $\underset{\longrightarrow}{\lim}\simeq x$, est immédiate.

La proposition précédente ne dit pas si on peut prendre pour $J$ un ens. ordonné. Dans ce cas, la condition c) signifie que pour \unsure{tte} partie $J'$ de $J$ \unsure{grde|grd} devant $\pi_0$ et de cardinal $\leq\pi$, $J$ admet une borne supérieure $i_J$ dans $J$, et la condition d) que l'on a alors

\[
x_{i_J}\xleftarrow{\sim}\underset{i\in J'}{\underset{\longleftarrow}{\lim}}\,x_i.
\]

L'idée naturelle, pour construire un système inductif $\pi_0$-$\pi$-adapté $\inserted{(x_j)_{j\in\unsure{J|I}}}$, avec $I$ ordonné, en partant de celui qu'on a construit, avec $J=M_\pi/x$, est de prendre pour $I$ l'ens., ordonné par inclusion, des sous-catégories $\inserted{J'}$ de \deleted{$\unsure{J_0}$}$\inserted{J_0}\inserted{\text{ (sous-catégorie }\deleted{\text{\unsure{pleine}}}\text{ pleine petite de }J\text{ équivalente à }J\text{)}}$, qui sont de cardinal $\leq\pi$ et \unsure{grdes|grde} devant $\pi_0$. \inserted{(Comparer (4.7.3.3) \unsure{page|pge} 165)} Cet ens. ordonné est bel et bien \inserted{(petit et)} \uline{grd devant} $\pi$, si on suppose que $\pi^{\pi_0}=\pi$ (cf dém. 4.7.5 p.~168--\unsure{17|170})\unsure{--|.}

% archive sheet number: 389
\scanpage{391}\mspage{305}

\uuline{Lemme} 4.13.13. Soient $\pi_0,\pi>\pi_0$ des cardinaux infinis, \inserted{tels que $\pi^{\pi_0}=\pi$,} et $\deleted{\unsure{I}}J$ une \inserted{petite} catégorie, grde devant $\pi$. Soit $\deleted{\unsure{J}}I$ l'ens. des sous-catégories $J'$ de $J$ \deleted{telles que} qui sont grdes devant $\pi_0$, et de cardinal $\leq\pi$. Alors $I$ est $\pi_0$-$\pi$-adaptée.

Il reste à voir que pour \unsure{tt} ens. ordonné $A$ grd devant $\pi_0$ et de cardinal $\leq\pi$, si $\unsure{i|r}\mapsto J_r$ est une famille croissante d'éléments de $I$ indexée par $A$, alors

\[
J\overset{\text{\unsure{déf|def}}}{=}\bigcup_{r\in A}J_r=\underset{r\in A}{\underset{\longrightarrow}{\lim}}\,J_r
\]

est dans $I$, i.e. est une sous-catégorie \inserted{de $I$} de cardinal $\leq\pi$ (ce qui est clair) et \uline{grde devant} $\pi_0$. Cela va résulter du

\uline{Lemme} 4.13.13.1 -- Soit $(J_r)_{r\in A}$ un système inductif de catégories grdes devant $\pi_0$, indexé par un ens. ordonné $A$ grd devant $\pi_0$. Alors $J=\underset{\unsure{r\in A}}{\underset{\longrightarrow}{\lim}}\,J_r$ est grd devant $\pi_0$.

% archive sheet number: 390
\scanpage{392}\mspage{306}

Il faut voir que si $B$ est un ens. de cardinal $\leq\pi_0$, alors le foncteur diagonal

\[
\delta:J\xrightarrow{\deleted{\unsure{\delta}}}J^B \tag*{$(*)$}
\]

est cofinal. Or comme $A$ est grd devant $\pi_0$, \inserted{$\mathrm{card}\,B\leq\pi_0$,} les produits \unsure{indexés} par $B$ commutent aux $\underset{\longrightarrow}{\lim}$ de type $A$ (SGA 4 I 9.8 appliqué à $\mathrm{Cat}$ et à $\pi_0$). Donc la flèche $(*)$ \deleted{\unsure{est la}} s'interprète comme limite inductive des flèches similaires

\[
\delta_r:J_r\longrightarrow J_r^B, \tag*{$(**)$}
\]

lesquelles sont cofinales, donc $\delta$ est cofinal comme limite inductive de foncteurs cofinaux.

\uline{Lemme} 4.13.13.2. Sous les hypothèses de 4.13.13, soient $M$ une catégorie\inserted{ satisfaisant \unsure{à} $\unsure{L^{\rightarrow}_{\pi_0}|L_{\pi_0}}$}, et $(x_j)_{j\in J}$ \inserted{(i.e. $J\longrightarrow M$)} un système inductif \deleted{$J\longrightarrow M$} dans $M$, indexé par $J$. Pour tout $J'\in I$, posons

\[
x_{J'}=\underset{i\in J'}{\underset{\longrightarrow}{\lim}}\,x_j
\]

% archive sheet number: 391
\scanpage{393}\mspage{307}

(qui existe dans $M$, puisque $M$ satisfait à $L_{\pi_0}$). Alors \deleted{\unsure{les}} $(x_{J'})_{J'\in I}$ forme un système inductif dans $M$, et $I\longrightarrow M$ commute aux limites inductives grdes devant $\pi_0$ et de cardinal $\leq\pi$.

Si $M$ est $\pi_0$-accessible, et les $x_j$ ($j\in\mathrm{Ob}\,\unsure{J}$) sont dans $M_\pi$, alors $(x_{J'})_{J'\in I}$ est un système inductif $\pi_0$-$\pi$-adapté dans $M$.

Tout ceci est \unsure{ess$^{\mathrm{t}}$.} trivial, et donne

\uline{Corollaire} 4.13.14 (de 4.13.1\unsure{2|7|3}) : sous les conditions de 4.13.12, si $\pi^{\pi_0}=\pi$, on peut pour \unsure{tt} $x\in\mathrm{Ob}\,M$ trouver une représentation $x\simeq\underset{j\in I}{\underset{\longrightarrow}{\lim}}\,x_j$, où $(x_i)_{i\in I}$ est un système inductif $\pi_0$-$\pi$-adapté avec $I$ \uline{ens. ordonné}.

\deleted{\unsure{Et}} De plus, si $\pi'>\pi$ et si $x\in M_{\pi'}$, on peut prendre $I$ tel qu'on ait

\[
\mathrm{card}\,I\leq{\pi'}^{\pi}.
\]

Ce dernier complément est immédiat,

% archive sheet number: 392
\scanpage{394}\mspage{308}

par les majorations habituelles.

% slanted
\margin{Cf. version plus claire et plus systématique du \uuline{th. 4.13.15} dans 4.13.22 p.~342.}

\deleted{\uline{Proposition}} \inserted{\uline{Théorème}} 4.13.15. Soient $M$ une catégorie $\pi_0$-accessible, et $\pi>\pi'_0\geq\pi_0$ tels que $\pi^{\pi_0^{\unsure{*|\prime}}}=\pi$*. Soit $\underline S$ une sous-catégorie strictement pleine de $M$, et soit $\underline S(\pi)=\underline S\cap M_\pi$. Les conditions suivantes sont équivalentes :

% slanted
\margin{* On peut se borner à \unsure{supp.|supposer} $\pi$ utile pour $M$ (\unsure{càd|et} $\pi^{\pi_0}=\pi$), alors on trouve que a), a', b, c sont équivalentes \unsure{et} \unsure{impliquent} c'), mais \unsure{on|ne} \unsure{sont|sait} \unsure{peut-être} \unsure{toujours} \unsure{pas} équivalentes à c'.}

\textcircled{a} $\underline S$ est $\pi$-accessible dans $M$ (4.13.3), \deleted{\unsure{et} $\underline S$} et $\underline S$ est stable dans $M$ par limites inductives \deleted{de cardinal} \inserted{grdes} devant $\pi'_0$.

\deleted{\unsure{b}}a') $\underline S$ est $\pi$-accessible dans $M$, et $\underline S(\pi)$ est stable dans $M_\pi$ par limites inductives grdes devant $\pi'_0$ et de cardinal $\leq\pi$*.

% slanted
\margin{donnons un nom à cette dernière condition, soit $L(\pi'_0,\pi)$ [pour $\underline S(\pi)$ dans $M_\pi$].}

b) Pour tout objet $x$ de $M$, \inserted{on a ceci :} la sous-catégorie \inserted{str.} pleine $\underline S(\pi)_{/x}$ de $\inserted{J=}M_{\pi}\unsure{\ell|<}/x$ est cofinale dans $M_\pi/x$ sss $x\in\mathrm{Ob}\,\underline S$. \inserted{De plus, $\underline S(\pi)$ est stable dans $M_\pi$ par $\underset{\longrightarrow}{\lim}$ grdes devant $\pi'_0$ et de cardinal $\leq\pi$.}

\deleted{\unsure{d}}\textcircled{c} Pour \unsure{tt} objet $x$ de $M$ et pour \unsure{tte} représentation $x\simeq\underset{i\in I}{\underset{\longrightarrow}{\lim}}\,x_i$, où $\deleted{\unsure{I}}(x_i)_{i\in I}$ est un système inductif $\pi'_0$-$\pi$-adapté (4.13.11), on a ceci : $\deleted{x\in}$ la sous-catégorie \inserted{str.} pleine $I_{\underline S}$ \inserted{de $I$} formée des $i\in\mathrm{Ob}\,I$ tels que $x_i\in\mathrm{Ob}\,\underline S$ est cofinale dans $I$ sss $x\in\mathrm{Ob}\,\underline S$. De plus, $\underline S(\pi)$ est stable \unsure{en|dans} $M_\pi$ par $\underset{\longrightarrow}{\lim}$ grdes devant $\pi'_0$ et de cardinal $\leq\pi$.

% archive sheet number: 393
\scanpage{395}\mspage{309}

\deleted{\unsure{d|c}}\unsure{c'|d'}) Comme d), avec $I$ ens. ordonné.

\uline{Corollaire} 4.13.16. Les hypothèses \inserted{sur $M,\pi_0,\underline S$} étant celles de 4.13.15, pour qu'il existe des cardinaux $\pi>\pi'_0\geq\pi_0$ satisfaisant les conditions $\pi^{\pi'_0}=\pi$ \inserted{(a fortiori $\pi^{\pi_0}=\pi$)} et les conditions équivalentes a) à d') de 4.13.15, il (f et) s que $\underline S$ soit $\deleted{\pi}$accessible dans $M$. De façon plus précise, \deleted{\unsure{En}} si $\underline S$ est accessible, dans $M$, soit $\pi'_0\geq\pi_0$ un cardinal tel que $\underline S$ soit $\pi'_0$-accessible dans $M$, et soit $\pi$ un cardinal \inserted{$>\pi'_0$} tel que $\pi^{\pi'_0}=\pi$. Alors $(\underline S,\pi'_0,\pi)$ satisfait aux conditions équivalentes de 4.13.15.

En effet, le ``il faut'' est évident sur la condition a) de 4.13.15. Pour le ``il suffit'' on \unsure{note|voit} que $\underline S$ $\pi'_0$-accessible dans $M$ implique que $\underline S$ est stable dans $M$ par $\underset{\longrightarrow}{\lim}$ \unsure{grdes|grde} devant $\pi'_0$. Donc pour vérifier a), il suffit de vérifier que $\underline S$ est $\pi$-accessible dans $M$, Or comme $\underline S$ est

% archive sheet number: 394
\scanpage{396}\mspage{310}

$\pi'_0$-accessible et que $\pi>\pi'_0$ satisfait $\unsure{\pi_0|\pi}^{\pi'_0}=\pi$, on en conclut que $\underline S$ est $\underline S_\pi$-accessible, i.e. $\underline S\xleftarrow{\approx}\mathrm{Ind}_\pi(\underline S_\pi)$. Donc en vertu de 4.13.3, il reste à voir que $\underline S_\pi\subset M_\pi$.

\deleted{Or \inserted{comme $\pi=\pi^{\pi'_0}$,} $\underline S_\pi$ est formé des \inserted{objets de $\underline S$ \unsure{isom.} à des} $\underset{\longrightarrow}{\lim}$ $x_i$, où $I$ \unsure{qui}} ou ce qui revient au même, que $\underline S_\pi\subset\underline S(\pi)\overset{\text{\unsure{def|déf}}}{=}\underline S\cap M_\pi$. (En fait, on aura $=$ égalité, car il est \deleted{\unsure{immédiat}} \inserted{clair} que $\underline S(\pi)\subset\underline S_\pi\;\ldots$).

Or $\underline S(\pi)$ est une sous-catégorie de $\underline S$ stable par $\underset{\longrightarrow}{\lim}$ grdes devant $\pi'_0$ et de cardinal $\leq\pi$ (puisque $\underline S$ et $M_\pi$ sont stables par ces limites) \uline{dans $M$}, et que ces limites dans $\underline S$ et dans \uline{$M$} sont \unsure{itou|idem}), et $\underline S(\pi)$ contient $\underline S_{\pi'_0}$. Donc $\underline S(\pi)$ contient $\underline S_\pi$, en vertu de 4.7.1 (page 161), implication c)$\Longrightarrow$b', appliqué à $\underline S,\pi'_0,\pi$ \unsure{ald|a.l.d.} $M,\pi_0,\pi$. Cela prouve \unsure{bien|donc} le corollaire 4.13.16.

\uline{Démonstration de} 4.13.15. Voici les implications $\pm$ tautologiques

\diagram{\begin{tikzcd}[ampersand replacement=\&,column sep=small,row sep=0.4em]
\& \& \& \& \deleted{\unsure{d}}c'\\
(*)\qquad \& a)\arrow[r,Rightarrow] \& \deleted{\unsure{b}}a'\arrow[r,Rightarrow,"\unsure{\cdot}"] \& \deleted{\unsure{d}}c\arrow[ru,Rightarrow]\arrow[rd,Rightarrow] \& \\
\& \& \& \& b,
\end{tikzcd}}

% archive sheet number: 395
\scanpage{397}\mspage{311}

que je vais expliciter un peu. On a $a\Rightarrow\deleted{\unsure{b}}a'$ puisque \deleted{$\underline S$} si a est satisfaite, $\underline S$ \uline{et} $M_\pi$ sont stables \inserted{\uline{dans} $M$} par $\underset{\longrightarrow}{\lim}$ grdes devant $\pi'_0$ et de cardinal $\leq\pi$, donc aussi leur intersection, $\underline S(\pi)$, qui est donc aussi stable dans $M_\pi$ par lesdites limites. \deleted{\unsure{L|a}} L'implication $\deleted{\unsure{d}\Rightarrow\unsure{c}},$ $c\Rightarrow c'$ \deleted{\unsure{est donc}} est une tautologie, et $c\Rightarrow b$ aussi, compte tenu de 4.13.12 qui nous dit que le syst. inductif canonique $(x_i)_{i\in I}$, avec $I=M_\pi/x$, de limite $x$, est $\pi_0$-$\pi$-adapté

(NB Comme $\pi^{\pi'_0}=\pi$, \inserted{et $\pi'_0\geq\pi_0$} on a aussi $\pi^{\pi_0}=\pi$, donc $\pi$ est utile pour $M$, \deleted{\unsure{)}} de sorte que 4.13.12 s'applique.

Reste à prouver la ``$\pm$ tautologie'' $\deleted{\unsure{b}}a'\Rightarrow\deleted{\unsure{d}}c$. \deleted{\illegible{1}} \inserted{Soit} donc une représentation $\pi_0$-$\pi$-adaptée

\[
x\simeq\underset{i\in I}{\underset{\longrightarrow}{\lim}}\,x_i,
\]

prouvons que

\[
x\in\mathrm{Ob}\,\underline S\Longleftrightarrow I_{\underline S}\ \text{cofinal dans }I.
\]

% slanted
\margin{On n'utilise, \inserted{ici, pour $\Leftarrow$,} de l'hypothèse $(x_i)_{i\in I}$ $\pi_0$-$\pi$-adapté, que le seul fait que $I$ soit grd devant $\pi$.}

L'implication $\Leftarrow\deleted{\unsure{>}}$ provient du fait que, $\underline S$ étant $\pi$-accessible dans $M$, est stable \inserted{dans $M$} par $\underset{\longrightarrow}{\lim}$ grdes devant $\pi$*. Reste à prouver

% archive sheet number: 396
\scanpage{398}\mspage{312}

que si $x\in\mathrm{Ob}\,\underline S$, alors $I_{\underline S}$ est cofinal dans $I$. Considérons \unsure{donc|alors} le diagramme de catégories

\diagram{\begin{tikzcd}[ampersand replacement=\&,column sep=large]
I \arrow[r,"\unsure{\varphi|\phi}"] \& J=M_\pi/x\\
I'=I_{\underline S} \arrow[u,hook] \arrow[r,"\varphi'"] \& J'=\underline S(\pi)/x,\arrow[u,hook,"\alpha"']
\end{tikzcd}}

c'est un carré cartésien ($I'=I_{\underline S}$ est l'image inverse par $\varphi:I\unsure{-|\longrightarrow}J$ de la sous-catégorie strictement pleine $J'$ de $J$). \deleted{\unsure{Comme}}

On a les propriétés suivantes

$\cdot$ 1) $I$ est grd devant $\pi$, et $\varphi$ cofinal ($\Longrightarrow J$ grd devant $\pi$)

$\cdot$ 2) $I$ stable par $\underset{\longrightarrow}{\lim}$ grdes devant $\pi'_0$ et de cardinal $\leq\pi$, et $\varphi$ commute à ces $\underset{\longrightarrow}{\lim}$.

\deleted{$\cdot$ 3) $J^{\unsure{4|\prime}}$ est stable par $\underset{\longrightarrow}{\lim}$ du type précédent, et $J'$ est}

% slanted
\margin{C'est ici qu'intervient l'hypothèse de stabilité \uline{faite dans b)}, \unsure{(}car d'après \unsure{b)|a')}, $\underline S$ est $\pi$-accessible dans $M$}

3) $J'$ est \deleted{stable dans \uline{cofinale} dans $J$, et} stable pour \deleted{$\unsure{\lim}$} $\underset{\longrightarrow}{\lim}$ \unsure{des types précédents|du type précédent} \inserted{et $\alpha$ y commute.}

4) $J'$ cofinal dans $J$.

On veut en conclure que $I'=\varphi^{-1}(J')$ --

% archive sheet number: 397
\scanpage{399}\mspage{313}

est cofinale dans $I$. Comme c'est une sous-catégorie pleine (et $=$ str. pleine), il suffit de prouver la propriété F1, i.e. que pour tout $i\in\mathrm{Ob}\,I$, existe $i'\in\mathrm{Ob}\,I'$ \deleted{tel que} et une flèche $i\longrightarrow i'$.

On va construire un $i'$ par construction transfinie, sur l'intervalle ordinal $I_{\alpha_\pi}$, en construisant un système inductif

\[
(i_\alpha)_{\alpha\in I_{\alpha_\pi}}\qquad i_\alpha\in\mathrm{Ob}\,I
\]

d'objets de $I$, et \deleted{une famille} \inserted{un système} \deleted{pour tout $\alpha\in I_{\alpha_\pi}$ un $j_\alpha\in J'$ et une flèche $\unsure{\varphi|\phi}(i_\alpha)\longrightarrow j_\alpha$.} \inserted{inductif}

\[
\inserted{(j_\alpha)_{\alpha\in I_{\alpha_\pi}},\qquad j_\alpha\in\mathrm{Ob}\,J',}
\]

\deleted{\unsure{(systif)|(respectif)}} flèches de transition notées

\[
\tau_{\beta,\alpha}:i_\alpha\unsure{-|\longrightarrow}i_\beta,\qquad
\tau'_{\beta,\alpha}:j_\alpha\unsure{-|\longrightarrow}j_\beta,
\]

enfin pour tout $\alpha\unsure{\text{dans}|\in}I_{\unsure{<|\alpha}\pi}$ des flèches

\diagram{\begin{tikzcd}[ampersand replacement=\&,column sep=small]
(*)\quad \& \varphi(i_\alpha)\arrow[r,"u_\alpha"]\arrow[rr,bend right=24,dashed,"{\varphi(\tau_{\alpha+1,\alpha})}"']\arrow[rrr,no head,bend right=33,start anchor=south,end anchor=south] \& j_\alpha\arrow[r,"v_\alpha"]\arrow[rr,bend left=28,dashed,"{\tau'_{\alpha+1,\alpha}}"] \& \varphi(i_{\alpha+1})\,[\arrow[r,dotted,"u_{\alpha+1}"] \& j_{\alpha+1}\,]
\end{tikzcd}}

dans $J$, telles que

\[
\left\{\begin{aligned}
v_\alpha u_\alpha&=\varphi(\tau_{\alpha+1,\alpha})\unsure{=|,}\\
u_{\alpha+1}v_\alpha&=\tau'_{\alpha+1,\alpha},
\end{aligned}\right.\tag*{$(**)$}
\]

% slanted
\margin{J'aurais \unsure{(aussi)} pu \unsure{prendre}}

% archive sheet number: 398
\scanpage{400}\mspage{314}

Montrons comment la construction d'un tel système d'objets et de flèches, \inserted{avec $(***)$ $i_0=i$} va entraîner \deleted{la \unsure{condition|conclusion} voulue} \inserted{l'existence de $i'\in\mathrm{Ob}\,I'$ et de $i\longrightarrow i'$.} \deleted{que $I'$ est cofinal dans $I$.} En effet, comme $I_{\alpha_\pi}$ est grd devant $\pi'_0$ (il est grd devant \unsure{tt|tout} cardinal $<\pi$), \deleted{$\underset{\longrightarrow}{\lim}\,i_\alpha$} il existe

\[
i'\overset{\text{\unsure{def|déf}}}{=}\underset{i\in I_{\alpha_\pi}}{\underset{\longrightarrow}{\lim}}\,i_\alpha
\]

dans $I$, et on a une flèche de transition

\[
i=i_0\longrightarrow i',
\]

donc il reste à prouver que $i'\in\mathrm{Ob}\,I'$, i.e. $\unsure{\varphi|\phi}(i')\in\mathrm{Ob}\,J'$. Or $\varphi$ commutant aux $\underset{\longrightarrow}{\lim}$ de type $I_{\unsure{<|\alpha}\pi}$, on aura

\[
\varphi(i)=\underset{\longrightarrow}{\lim}\ \varphi(i_\alpha).
\]

Il nous faut encore des compatibilités

\diagram{\begin{tikzcd}[ampersand replacement=\&,column sep=small,row sep=1.6em]
(****) \quad \& \varphi(i_\alpha)\arrow[r,dash,"u_\alpha"]\arrow[d,"{\varphi(\tau_{\beta,\alpha})}"'] \& j_\alpha\arrow[d,"\tau'_{\beta\alpha}"] \& \text{et} \& j_\alpha\arrow[r,"v_\alpha"]\arrow[d,"{\tau'_{\beta,\alpha}}"] \& \varphi(i_{\alpha+1})\arrow[d,"{\varphi(\tau_{\beta+1,\alpha+1})}"]\\
\& \varphi(i_\beta)\arrow[r,"u_\beta"] \& j_\beta \& \& j_\beta\arrow[r,"v_\beta"] \& \varphi(i_{\beta+1}),
\end{tikzcd}}

oubliées \unsure{tantôt}.

% archive sheet number: 399
\scanpage{401}\mspage{315}

Les compatibilités $(****)$ impliquent qu'on a des hom. \deleted{\unsure{de}}

\diagram{\begin{tikzcd}[ampersand replacement=\&,column sep=large]
\unsure{\varphi|\phi}(i)\arrow[r,shift left=0.65ex,"u"] \& j\arrow[l,shift left=0.65ex,"v"] \& {,\ \text{où }\ j\overset{\text{\unsure{def|déf}}}{=}\underset{\alpha\in I_{\alpha_\pi}}{\underset{\longrightarrow}{\lim}}\,j_\alpha}
\end{tikzcd}}

\inserted{(\unsure{existe|existant} \deleted{\unsure{la}}\unsure{déjà|la|là} \uline{dans $J'$} donc dans $J$\unsure{)|)'}},

et les compatibilités $(**)$ impliquent que

\[
\deleted{\unsure{u}\ \text{et}\ \unsure{v}\ \text{\unsure{son}}}\qquad
vu=\mathrm{id}_{\varphi(i)}\deleted{\unsure{\varphi(i)}},\quad vu=\mathrm{id}_j,
\]

donc que $u,v$ sont des iso. inverses l'un de l'autre. \deleted{Or la propriété 4) pour $J$ implique que $j\in J'$,} Comme $J'$ est une sous-catégorie \uline{str.} pleine de $J$, il s'ensuit bien que $\varphi(i)\in J'$, i.e. $i\in I'$, qed.

\deleted{\unsure{Donc}} \unsure{Reste|Rest} \unsure{maint$^{\mathrm{t}}$|maintenant} à construire les

\[
\underset{\mathrm{Ob}\,I}{\underset{\unsure{\mathord{\tikz[baseline=(c.base)]{\node[rotate=-90,inner sep=0pt] (c) {$\scriptstyle\in$};}}}}{i_\alpha}},\tau_{\beta\alpha},\quad
\underset{\mathrm{Ob}\,J'}{\underset{\unsure{\mathord{\tikz[baseline=(c.base)]{\node[rotate=-90,inner sep=0pt] (c) {$\scriptstyle\in$};}}}}{j_\alpha}},\tau'_{\beta\alpha},\quad u_\alpha,v_\alpha
\]

satisfaisant aux contraintes \deleted{$\unsure{(*)}$} $(**)\,(***)\,(****)$. On va faire la construction en trois étapes.

\textcircled{1} Départ :

\[
i_0=\unsure{0|i},
\]

Comme $J'$ est cofinal dans $J$, existe $j_0\in\mathrm{Ob}\,J'$ et $\varphi(i_0)\xrightarrow{u_0}j_0$.

% archive sheet number: 400
\scanpage{402}\mspage{316}

\textcircled{2} Supposons que les $i_\beta,j_\beta$ soient construits pour $\beta<\alpha$, ainsi que les $\tau_{\beta',\beta}:\unsure{i_\beta|i'_\beta}\unsure{-|\longrightarrow}\unsure{i_{\beta'}|i'_{\beta'}}$, $\tau'_{\beta',\beta}:j_\beta\unsure{-|\longrightarrow}j_{\beta'}$ \inserted{(pour $\beta,\beta'<\alpha$)}, \inserted{\unsure{outre|enfin}} les $u_\beta:\unsure{\varphi|\phi}(i_\beta)\longrightarrow j_\beta$, $v_\beta:j_\beta\longrightarrow\varphi(i_{\beta+1})$ \deleted{\unsure{cet} \unsure{$v_\beta$} \unsure{pour}} pour $\beta<\alpha$ resp. pour $\beta+1<\alpha$, de façon à satisfaire aux contraintes $(**)\,(****)$ \unsure{i.e.|les}

\deleted{\unsure{faut} construire 1) $i_\alpha\deleted{\unsure{\in J}}\inserted{\in\mathrm{Ob}\,I}$ et des $i_\beta\xrightarrow{\tau_{\unsure{\beta,\alpha|\alpha,\beta}}}i_\alpha$ ($\beta<\alpha$), compatibles avec les $\tau_{\beta',\beta}$ ($\beta<\beta'<\alpha$) 2) $j_\alpha\in\mathrm{Ob}\,J'$ et les $j_\beta\xrightarrow{\tau'_{\beta\alpha}}j_\alpha$ ($\beta<\alpha$) compatibles avec les $\tau'_{\beta',\beta}$ ($\beta<\beta'<\alpha$) 3) $u_\alpha:\varphi(i_\alpha)$}

commutativités des

\diagram{\(\left\{\begin{array}{l}
\left.
\begin{gathered}
\begin{tikzcd}[ampersand replacement=\&,column sep=small]
\text{\textcircled{1}}\ \varphi(i_\beta)\arrow[r,"u_\beta"]\arrow[rr,bend right=23,"{\varphi(\tau_{\beta+1,\beta})}"'] \& j_\beta\arrow[r,"v_{\beta\deleted{\unsure{+1}}}"] \& \varphi(i_{\beta+1})
\end{tikzcd}\\[1.1em]
\begin{tikzcd}[ampersand replacement=\&,column sep=small]
\text{\textcircled{2}}\ \deleted{\unsure{\varphi(i}}j_\beta\arrow[r,"v_{\beta\deleted{\unsure{+1}}}"]\arrow[rr,dashed,bend right=23,"{\tau'_{\beta+1,\beta}}"'] \& \varphi(i_{\beta+1})\arrow[r,"u_{\beta+1}"] \& j_{\beta+1}
\end{tikzcd}
\end{gathered}\right\}\ \text{si }\underline{\beta+1<\alpha}\\[1em]
\begin{tikzcd}[ampersand replacement=\&,column sep=small,row sep=1.2em]
\text{\textcircled{3}}\quad \varphi(i_\beta)\arrow[r,"u_\beta"]\arrow[d,"{\varphi(\tau_{\beta',\beta})}"'] \& j_\beta\arrow[d,"{\tau'_{\beta',\beta}}"] \& \text{si }\beta<\beta'<\alpha\\
\varphi(i_{\beta'})\arrow[r,"u_{\beta'}"] \& j_{\beta'} \&
\end{tikzcd}\\[0.7em]
\begin{tikzcd}[ampersand replacement=\&,column sep=small,row sep=1.2em]
\text{\textcircled{4}}\quad j_\beta\arrow[r,"v_\beta"]\arrow[d,"{\tau'_{\beta',\beta}}"'] \& \varphi(i_{\beta+1})\arrow[d,"{\varphi(\tau_{\beta'+1,\beta+1})}"] \& \text{si }\beta<\beta'\\
j_{\beta'}\arrow[r,"v_{\beta'}"] \& \varphi(i_{\beta'+1}) \& \text{et }\beta'+1<\alpha
\end{tikzcd}
\end{array}\right.\)}

% archive sheet number: 401
\scanpage{403}\mspage{317}

il faut construire

1) $i_\alpha\in\mathrm{Ob}\,I$ et un hom $(i_\beta)_{\beta\in I_\alpha}\longrightarrow i_\alpha$ i.e. des $\tau_{\deleted{\unsure{\beta'}}\alpha,\beta}:i_\beta\unsure{-|\longrightarrow}i_\alpha$, \inserted{$(\beta\in I_\alpha)$}, compatibles aux transitions $\tau_{\beta',\beta}$ pour $\beta<\beta'<\alpha$

2) $j_\alpha\in\mathrm{Ob}\,J'$, et un hom $(j_\beta)_{\beta\in I_\alpha}\unsure{-|\longrightarrow}j_\alpha$, i.e des $\tau'_{\alpha,\beta}:i_\beta\unsure{-|\longrightarrow}i_\alpha$, $(\beta\in I_\alpha)$, compatibles aux transitions $\tau'_{\beta',\beta}$ pour $\beta<\beta'<\alpha$ --

3) \deleted{\unsure{Une}} flèche

\[
u_\alpha:\unsure{\varphi|\phi}(i_\alpha)\unsure{-|\longrightarrow}j_\alpha
\]

4) Si $\alpha$ n'est pas ordinal limite, i.e. $\alpha=\unsure{\alpha'_0|\alpha_0}+1$, une flèche

\[
v_{\alpha_0}:\deleted{\unsure{\varphi}}j_{\alpha_0}\longrightarrow\varphi(i_{\alpha_0+1})=\varphi(i_\alpha)\,;
\]

\deleted{de façon} et ceci de façon à satisfaire aux contraintes \textcircled{1}\textcircled{2}\textcircled{3}\textcircled{4} quand on y suppose $\beta+1\leq\alpha$ (ald $\beta+1<\alpha$) dans \textcircled{1} et \textcircled{2}, et $\beta'\leq\alpha$ (ald $\beta'<\alpha$) dans \textcircled{3}, et $\beta'+1\leq\alpha$ (ald $\beta'+1<\alpha$) dans \textcircled{4}. Donc \inserted{pour} les contraintes $1^\circ,2^\circ,4^\circ$, \deleted{\unsure{vient à être}\ \unsure{en}} il n'y a en tenir compte \uline{que} si $\alpha$ n'est pas ordinal limite, donc $\alpha=\alpha_0+1$, et en y faisant

\[
\beta'=\alpha_0.
\]

Quant à la contrainte $3^\circ$, il faut en

% archive sheet number: 402
\scanpage{404}\mspage{318}

tenir compte pour $\beta'=\alpha$. Donc au total, il faut exiger

\diagram{\(\displaystyle (3_{\alpha})(\beta)\quad\begin{tikzcd}[ampersand replacement=\&]
\unsure{\varphi|\phi}(i_{\beta}) \arrow[r,dash,"{u_{\beta}}"] \arrow[d,"{\varphi(\mathord{\tikz[baseline=(c.base)]{\node[ellipse,draw,inner sep=1pt] (c) {$\scriptstyle \tau_{\alpha,\beta}$};}})}"'] \& j_{\beta} \arrow[d,"{\mathord{\tikz[baseline=(c.base)]{\node[ellipse,draw,inner sep=1pt] (c) {$\scriptstyle \tau'_{\alpha,\beta}$};}}}"]\\
\varphi(i_{\alpha}) \arrow[r,"{\mathord{\tikz[baseline=(c.base)]{\node[ellipse,draw,inner sep=1pt] (c) {$\scriptstyle u_{\alpha}$};}}}"] \& j_{\alpha}
\end{tikzcd}\quad\begin{gathered}\text{commutatif}\\\text{pour }\text{\uline{$\beta$}}<\alpha,\end{gathered}\)}

% slanted
\margin{\uline{NB} Il suffit de l'exiger pour $\beta=\alpha_0$, dans le cas $\alpha=\alpha_0+1$, cf ci-dessous p.~323.}

(par restriction sur $\alpha$), et si $\alpha$ \inserted{$=\alpha_0+1$} est ordinal non limite, il faut exiger de \inserted{plus}

\diagram{\(\displaystyle (4_{\alpha})(\beta)\quad\begin{tikzcd}[ampersand replacement=\&]
j_{\beta} \arrow[r,"{v_{\beta}}"] \arrow[d,"{\deleted{\unsure{\varphi}}\;\tau'_{\alpha_0,\beta}}"'] \& \varphi(i_{\beta+1}) \arrow[d,"{\varphi(\mathord{\tikz[baseline=(c.base)]{\node[ellipse,draw,inner sep=1pt] (c) {$\scriptstyle \tau_{\alpha,\beta+1}$};}})}"]\\
j_{\alpha_0} \arrow[r,"{\mathord{\tikz[baseline=(c.base)]{\node[ellipse,draw,inner sep=1pt] (c) {$\scriptstyle v_{\alpha_0}$};}}}"] \& \varphi(i_{\alpha_0+1})=\varphi(i_{\alpha})
\end{tikzcd}\quad \text{\uline{$\beta$}}<\alpha_0\)}

\diagram{\(\displaystyle (1_{\alpha})\quad\begin{tikzcd}[ampersand replacement=\&,column sep=large]
\varphi(i_{\deleted{\unsure{\beta}}\alpha_0}) \arrow[r,"{u_{\alpha_0}}"] \arrow[rr,bend right=28,"{\varphi(\mathord{\tikz[baseline=(c.base)]{\node[ellipse,draw,inner sep=1pt] (c) {$\scriptstyle \tau_{\alpha,\alpha_0}$};}})}"'] \& j_{\alpha_0} \arrow[r,"{\mathord{\tikz[baseline=(c.base)]{\node[ellipse,draw,inner sep=1pt] (c) {$\scriptstyle v_{\alpha_0}$};}}}"] \& \varphi(i_{\alpha_0+1})=\varphi(i_{\alpha})
\end{tikzcd}\)}

\diagram{\(\displaystyle (2_{\alpha})\quad\begin{tikzcd}[ampersand replacement=\&]
j_{\alpha_0} \arrow[r,"{\mathord{\tikz[baseline=(c.base)]{\node[ellipse,draw,inner sep=1pt] (c) {$\scriptstyle v_{\alpha_0}$};}}}"] \arrow[rr,bend right=26,"{\deleted{\unsure{\varphi(\tau|\varphi}}\;\mathord{\tikz[baseline=(c.base)]{\node[ellipse,draw,inner sep=1pt] (c) {$\scriptstyle \tau'_{\alpha,\alpha_0}$};}}}"'] \& \varphi(i_{\alpha_0+1})=\varphi(i_{\alpha}) \arrow[r,"{\mathord{\tikz[baseline=(c.base)]{\node[ellipse,draw,inner sep=1pt] (c) {$\scriptstyle u_{\alpha}$};}}}"] \& j_{\alpha}
\end{tikzcd}\quad .\)}

Dans \unsure{ces|les} compatibilités, j'ai encerclé les flèches $\tau_{\alpha,\beta},\tau'_{\alpha,\beta}\quad(\beta<\alpha),\;[\tau_{\alpha,\beta+1}\;(\beta+1<\alpha)]$ $v_{\alpha_0},\deleted{\unsure{,|;}}\,[\tau_{\alpha,\alpha_0}],\,u_{\alpha},\,[\tau'_{\alpha,\alpha_0}]$, qui ont été introduites \deleted{au cours} \deleted{\unsure{des}}\unsure{dans} les données 1) 2) 3) 4) page 317.

% archive sheet number: 403
